% 第17章17.5 Auction算法算例：成本 C=[[1,2],[1,4]]（竞标者：甲、乙；物品：任务1、任务2），epsilon=1/4，最小成本形式。
% 比较值 q_ij = c_ij + p_j（出价时的价格）。
% (a) 甲：q=(1,2)，选任务1，涨价 2-1+1/4=5/4，p=(5/4,0)。
% (b) 乙：q=(9/4,4)，选任务1，涨价 4-9/4+1/4=2，p=(13/4,0)；甲失去任务1。
% (c) 甲：q=(17/4,2)，选任务2，涨价 17/4-2+1/4=5/2，p=(13/4,5/2)。结束：甲-任务2，乙-任务1，原成本3。
% 数值由 scratchpad 中的 verify.py 复算。
\begin{tikzpicture}[line width=.8pt,font=\fontsize{8}{10}\selectfont,
  rw/.style={rectangle,rounded corners=1.8pt,draw=FigureStroke,fill=FigureBlueFill,line width=.8pt,minimum width=6.5mm,minimum height=6mm,inner sep=1pt,text=FigureInk},
  cl/.style={rw,fill=FigureGreenFill,minimum width=10mm},
  cmp/.style={draw=FigureMuted,line width=.7pt,dashed},
  hold/.style={draw=FigureBlue,line width=2pt},
  bid/.style={draw=FigureGreen,line width=2pt},
  lost/.style={draw=FigureRed,line width=.9pt,dash pattern=on 1.5pt off 1.5pt},
  lab/.style={text=FigureInk,inner sep=1.2pt,fill=white},
  pr/.style={text=FigureMuted,font=\FigureNoteFont\fontsize{7.5}{9}\selectfont}]
  \foreach \dx/\k in {0/A,3.5/B,7.0/C}{
    \node[rw] (r1\k) at (\dx,.65) {甲};
    \node[rw] (r2\k) at (\dx,-.65) {乙};
    \node[cl] (c1\k) at (\dx+2.1,.65) {任务$1$};
    \node[cl] (c2\k) at (\dx+2.1,-.65) {任务$2$};
  }
  % (a) 甲出价
  \draw[bid] (r1A)--node[lab,pos=.3]{$1$} (c1A);
  \draw[cmp] (r1A)--node[lab,pos=.3]{$2$} (c2A);
  \node[pr,above=1pt] at (c1A.north) {$p_1{:}\,0\to5/4$};
  \node[pr,below=1pt] at (c2A.south) {$p_2=0$};
  \node at (1.05,-1.9) {(a) 甲出价};
  % (b) 乙出价，甲失去任务1
  \draw[lost] (r1B)--(c1B);
  \draw[bid] (r2B)--node[lab,pos=.3]{$9/4$} (c1B);
  \draw[cmp] (r2B)--node[lab,pos=.3]{$4$} (c2B);
  \node[pr,above=1pt] at (c1B.north) {$p_1{:}\,5/4\to13/4$};
  \node[pr,below=1pt] at (c2B.south) {$p_2=0$};
  \node at (4.55,-1.9) {(b) 乙出价，甲出局};
  % (c) 甲改选任务2
  \draw[hold] (r2C)--(c1C);
  \draw[cmp] (r1C)--node[lab,pos=.3]{$17/4$} (c1C);
  \draw[bid] (r1C)--node[lab,pos=.3]{$2$} (c2C);
  \node[pr,above=1pt] at (c1C.north) {$p_1=13/4$};
  \node[pr,below=1pt] at (c2C.south) {$p_2{:}\,0\to5/2$};
  \node at (8.05,-1.9) {(c) 甲改选任务$2$};
\end{tikzpicture}
